% =============================================================
% ลำดับการทำงานของสัญญาณนับคำซ้อน
% =============================================================

\begin{figure}[H]
\centering
\begin{tikzpicture}[
  bx/.style  = {draw=black!55, rounded corners=2pt, minimum width=3.3cm,
                minimum height=0.8cm, align=center, inner sep=2pt, font=\scriptsize},
  inp/.style = {bx, fill=black!6},
  op/.style  = {bx},
  mrg/.style = {bx, fill=black!20, very thick, draw=black, minimum width=4.4cm},
  fin/.style = {bx, fill=black!68, text=white, minimum width=4.4cm},
  note/.style= {draw=black!40, dashed, rounded corners=2pt, align=left,
                inner sep=5pt, font=\scriptsize, text width=4.4cm, fill=white},
  ar/.style  = {-{Stealth[length=4pt]}, thick, black!65}
]

\def\lx{-2.3} \def\rx{2.3}

% ---------- ข้อมูลเข้า ----------
\node[inp] (q0) at (\lx, 0)      {คำค้น\\(หลังผ่านขั้นแปลภาษา)};
\node[inp] (t0) at (\rx, 0)      {ชื่อสินค้าที่ค้นได้};

% ---------- ตัดเป็นคำ ----------
\node[op]  (q1) at (\lx,-1.35)   {ตัดเป็นคำ\\เฉพาะ a--z และ 0--9};
\node[op]  (t1) at (\rx,-1.35)   {ตัดเป็นคำ\\เฉพาะ a--z และ 0--9};

% ---------- ทิ้งคำ ----------
\node[op]  (q2) at (\lx,-2.75)   {ทิ้งคำที่ไม่ช่วยแยก};
\node[op]  (t2) at (\rx,-2.75)   {ทิ้งคำที่ไม่ช่วยแยก};

% ---------- เซต ----------
\node[op]  (q3) at (\lx,-4.1)    {เซตคำของคำค้น\\$Q$};
\node[op]  (t3) at (\rx,-4.1)    {เซตคำของชื่อสินค้า\\$T$};

% ---------- รวม ----------
\node[mrg] (m)  at (0,-5.5)      {หาคำที่อยู่ทั้งสองฝั่ง\\$|Q \cap T|$};
\node[fin] (f)  at (0,-6.9)      {คะแนน $= \dfrac{|Q \cap T|}{|Q|}$ \; มีค่า 0 ถึง 1};

\foreach \i/\j in {q0/q1, q1/q2, q2/q3, t0/t1, t1/t2, t2/t3} \draw[ar] (\i) -- (\j);
\draw[ar] (q3) |- (m.west);
\draw[ar] (t3) |- (m.east);
\draw[ar] (m)  -- (f);

% ---------- กฎที่ใช้ทิ้งคำ ----------
\node[note] (n) at (7.2,-2.75)
  {\textbf{กฎทิ้งคำ 2 ข้อ}\\[2pt]
   1. อยู่ในรายการคำที่ไม่ช่วยแยก 13 คำ\\
   \hspace*{0.9em}คำไวยากรณ์ และคำที่พบบ่อย\\
   \hspace*{0.9em}ในชื่อสินค้าจนไม่บอกความตรง\\[3pt]
   2. สั้นกว่า 3 ตัวอักษร\\
   \hspace*{0.9em}เพราะตรงกันโดยบังเอิญได้ง่าย};
\draw[ar, dashed, black!45] (n.west) -- (t2.east);

% ---------- หมายเหตุตัวหาร ----------
\node[note, text width=4.4cm] (d) at (7.2,-6.4)
  {\textbf{ตัวหารคือ $|Q|$}\\[2pt]
   นับจากจำนวนคำของ\textbf{คำค้น}\\
   ไม่ใช่ของชื่อสินค้า\\
   สินค้าที่ชื่อยาวจึงไม่ได้เปรียบ};
\draw[ar, dashed, black!45] (d.west) -- (f.east);

\end{tikzpicture}
\caption{ลำดับการทำงานของสัญญาณนับคำซ้อน}
\label{fig:wordcount}
\end{figure}
